%2multibyte Version: 5.50.0.2953 CodePage: 1252

\documentclass[notes=show]{beamer}
%%%%%%%%%%%%%%%%%%%%%%%%%%%%%%%%%%%%%%%%%%%%%%%%%%%%%%%%%%%%%%%%%%%%%%%%%%%%%%%%%%%%%%%%%%%%%%%%%%%%%%%%%%%%%%%%%%%%%%%%%%%%%%%%%%%%%%%%%%%%%%%%%%%%%%%%%%%%%%%%%%%%%%%%%%%%%%%%%%%%%%%%%%%%%%%%%%%%%%%%%%%%%%%%%%%%%%%%%%%%%%%%%%%%%%%%%%%%%%%%%%%%%%%%%%%%
\usepackage{amsfonts}
\usepackage{graphicx}
\usepackage{amssymb}
\usepackage{amsmath}
\usepackage{mathpazo}
\usepackage{hyperref}
\usepackage{multimedia}

\setcounter{MaxMatrixCols}{10}
%TCIDATA{OutputFilter=LATEX.DLL}
%TCIDATA{Version=5.50.0.2953}
%TCIDATA{Codepage=1252}
%TCIDATA{<META NAME="SaveForMode" CONTENT="2">}
%TCIDATA{BibliographyScheme=Manual}
%TCIDATA{Created=Wednesday, May 01, 2013 14:54:39}
%TCIDATA{LastRevised=Wednesday, October 19, 2022 18:02:47}
%TCIDATA{<META NAME="GraphicsSave" CONTENT="32">}
%TCIDATA{<META NAME="DocumentShell" CONTENT="Other Documents\SW\Slides - Beamer">}
%TCIDATA{Language=American English}
%TCIDATA{CSTFile=beamer.cst}
%TCIDATA{PageSetup=72,72,72,83,0}
%TCIDATA{Counters=arabic,1}
%TCIDATA{AllPages=
%H=36
%F=36
%}


\newenvironment{stepenumerate}{\begin{enumerate}[<+->]}{\end{enumerate}}
\newenvironment{stepitemize}{\begin{itemize}[<+->]}{\end{itemize} }
\newenvironment{stepenumeratewithalert}{\begin{enumerate}[<+-| alert@+>]}{\end{enumerate}}
\newenvironment{stepitemizewithalert}{\begin{itemize}[<+-| alert@+>]}{\end{itemize} }
\input{tcilatex}
\usetheme{JuanLesPins}
\usecolortheme{whale}

\begin{document}

\title{Econometrics Lecture 3: \\
Large Sample Properties}
\author{Katerina Petrova \ \ \ \ \ \ \ \ \ \ \ \ \ \ \ \ \ \ \ \ \ \ \ \ \ \
\ \ \ \ \ \ \ \ \ \ \ \ \ \ \ \ \ \ \ \ \ \ \ \ \ \ \ \ \ \ \ \ \ \ \ \ \ \
\ \ \ \ \ \ \ \ \ \ \ \ \ \ \ \ \ \ \ \ \ \ \ }
\date{}
\maketitle

\section{Formal Assumptions for the Linear Regression}

%TCIMACRO{\TeXButton{BeginFrame}{\begin{frame}}}%
%BeginExpansion
\begin{frame}%
%EndExpansion

\QTR{frametitle}{Assumptions for OLS}

\begin{enumerate}
\item The linear regression model in is correctly specified, and crucially,
linear in $\beta $ and $\varepsilon .$

\item The columns of $X$ are linearly independent, so that $rank(X^{\prime
}X)=rank(X)=k<n.$

\item $\varepsilon _{i}$ are i.i.d. variables and $X$ are exogenous
identically distributed regressors, such that $\mathbb{E}[\varepsilon
_{i}|X]=0$ and $\mathbb{E}[\varepsilon _{i}^{2}|X]=\sigma ^{2}.$

\item $p\lim_{n\rightarrow \infty }(\frac{1}{n}X^{\prime
}X)=p\lim_{n\rightarrow \infty }\frac{1}{n}\sum_{i=1}^{n}\mathbf{x}_{i}%
\mathbf{x}_{i}^{\prime }=\mathbb{E[}\mathbf{x}_{1}\mathbf{x}_{1}^{\prime
}]=Q,$ where $\left\Vert Q\right\Vert <\infty $ and $Q$ is positive definite
matrix.
\end{enumerate}

%TCIMACRO{\TeXButton{EndFrame}{\end{frame}}}%
%BeginExpansion
\end{frame}%
%EndExpansion

\section{Some useful results}

%TCIMACRO{\TeXButton{BeginFrame}{\begin{frame}}}%
%BeginExpansion
\begin{frame}%
%EndExpansion

\QTR{frametitle}{Some useful results}

Let $X_{1},...,X_{n}$ is a sequence of random variables and let $X$ be
another random variable. Let $F_{n}$ denote the cdf of $X_{n}$ and let $F$
denote the cdf of $X$.

\begin{definition}
$X_{n}$ converges to $X$ in distribution, written $X_{n}\rightarrow _{d}X,$
(you can also find it written $X_{n}\rightsquigarrow X$) if%
\begin{equation*}
\lim_{n\rightarrow \infty }F_{n}(t)=F(t)
\end{equation*}%
for all $t$ for which $F$ is continuous.
\end{definition}

\begin{itemize}
\item This can be extended to vector sequence $X_{1},...,X_{n}$ and vector
limit $X$ with $F_{n}$ and $F$ denoting cdf of vector-valued random
variables.
\end{itemize}

%TCIMACRO{\TeXButton{EndFrame}{\end{frame}}}%
%BeginExpansion
\end{frame}%
%EndExpansion

%TCIMACRO{\TeXButton{BeginFrame}{\begin{frame}}}%
%BeginExpansion
\begin{frame}%
%EndExpansion

\begin{definition}
$X_{n}$ converges to $X$ in probability, written $X_{n}\rightarrow _{p}X,$
if for every $\varepsilon >0$ (no matter how small),%
\begin{equation*}
P(|X_{n}-X|\geq \varepsilon )\rightarrow 0\text{ as }n\rightarrow \infty
\end{equation*}%
of equivalently, for any $\varepsilon ,\delta >0$, $\exists N(\varepsilon
,\delta )$, s.t. 
\begin{equation*}
\text{for each }n\geq N(\varepsilon ,\delta )\text{ \ \ \ \ }P(\left\vert
X_{n}-X\right\vert \geq \varepsilon )<\delta .
\end{equation*}%
Suppose that $X_{n}$ and $X$ are now random vectors in $\mathbb{R}^{k}.$ We
define $X_{n}\rightarrow _{p}X$ by element-wise convergence in probability: $%
X_{n,i}\rightarrow _{p}X_{i}$ as $n\rightarrow \infty $ for all $i\in
\left\{ 1,...,k\right\} .$ Equivalently, 
\begin{equation*}
\left\Vert X_{n}-X\right\Vert \rightarrow _{p}0,
\end{equation*}%
where $\left\Vert .\right\Vert $ denotes the Euclidian norm, $\left\Vert
x\right\Vert =\sqrt{x^{\prime }x}.$
\end{definition}

%TCIMACRO{\TeXButton{EndFrame}{\end{frame}}}%
%BeginExpansion
\end{frame}%
%EndExpansion

%TCIMACRO{\TeXButton{BeginFrame}{\begin{frame}}}%
%BeginExpansion
\begin{frame}%
%EndExpansion

\QTR{frametitle}{Some useful results}

\begin{enumerate}
\item $X_{n}\rightarrow _{p}X$ implies that $X_{n}\rightarrow _{d}X.$

\item If $X_{n}\rightarrow _{d}X$ and if $P(X=c)=1$ for some real number $c,$
then $X_{n}\rightarrow _{p}X.$

\item If $X_{n}\rightarrow _{p}X$ and $g(.)$ is a continuous function, then $%
g(X_{n})\rightarrow _{p}g(X)$ (Continuous Mapping Theorem)

\item If $X_{n}\rightarrow _{d}X$ and $g(.)$ is a continuous function, then $%
g(X_{n})\rightarrow _{d}g(X)$ (Continuous Mapping Theorem)

\item If $X_{n}\rightarrow _{d}X$ and $Y_{n}\rightarrow _{p}Y,$ then $%
X_{n}Y_{n}\rightarrow _{d}XY$ (Slutsky Theorem)

\item If $X_{n}\rightarrow _{d}X$ and $Y_{n}\rightarrow _{p}Y,$ then $%
X_{n}+Y_{n}\rightarrow _{d}X+Y$ (Slutsky Theorem)

\item If $X_{n}\rightarrow _{d}X$ and $Y_{n}\rightarrow _{p}Y,$ then $%
X_{n}/Y_{n}\rightarrow _{d}X/Y$ if $Y\neq 0$ (Slutsky Theorem)
\end{enumerate}

%TCIMACRO{\TeXButton{EndFrame}{\end{frame}}}%
%BeginExpansion
\end{frame}%
%EndExpansion

%TCIMACRO{\TeXButton{BeginFrame}{\begin{frame}}}%
%BeginExpansion
\begin{frame}%
%EndExpansion

\QTR{frametitle}{Weak Law of Large Numbers}

\begin{itemize}
\item Let $X_{1},...,X_{n}$ be an i.i.d. sample, let $\mu =\mathbb{E}[X_{1}]$
(note than $\mu =\mathbb{E}[X_{i}]$ is the same for all $i,$ so by
convention, we are often going to write $\mu =\mathbb{E}[X_{1}])$ and $%
\sigma ^{2}=V(X_{1}).$ Recall that the sample mean is defined $\overline{X}%
_{n}=\frac{1}{n}\dsum\nolimits_{i=1}^{n}X_{i}.$
\end{itemize}

\begin{theorem}[The Weak Law of Large Numbers]
If $X_{1},...,X_{n}$ are i.i.d. variables, and $\mathbb{E[}\left\vert
X_{1}\right\vert ]<\infty ,$ then $\frac{1}{n}\dsum\nolimits_{i=1}^{n}X_{i}%
\rightarrow _{p}\mu .$ Equivalently, $\left\vert \frac{1}{n}%
\dsum\nolimits_{i=1}^{n}X_{i}-\mathbb{E}[X_{1}]\right\vert \rightarrow
_{p}0. $
\end{theorem}

%TCIMACRO{\TeXButton{EndFrame}{\end{frame}}}%
%BeginExpansion
\end{frame}%
%EndExpansion

%TCIMACRO{\TeXButton{BeginFrame}{\begin{frame}}}%
%BeginExpansion
\begin{frame}%
%EndExpansion

\QTR{frametitle}{Weak Law of Large Numbers}

\begin{itemize}
\item Recall that 
\begin{equation*}
\mathbb{E}[\overline{X}_{n}]=\mathbb{E}[\frac{1}{n}\dsum%
\nolimits_{i=1}^{n}X_{i}]=\frac{1}{n}\dsum\nolimits_{i=1}^{n}\mathbb{E}%
[X_{i}]=\frac{n\mu }{n}=\mu
\end{equation*}%
and by independence of $X^{\prime }s$ 
\begin{equation*}
V(\overline{X}_{n})=V(\frac{1}{n}\dsum\nolimits_{i=1}^{n}X_{i})=\frac{1}{%
n^{2}}\dsum\nolimits_{i=1}^{n}V(X_{i})=\frac{n\sigma ^{2}}{n^{2}}=\frac{%
\sigma ^{2}}{n}
\end{equation*}

\item Another way to think about the Law of Large Numbers is that the random
variable $\overline{X}_{n}$ becomes more and more concentrated around a
single value, $\mu ,$ as $n$ gets large (note that its variance $\frac{%
\sigma ^{2}}{n}$ vanishes as $n\rightarrow \infty $).
\end{itemize}

%TCIMACRO{\TeXButton{EndFrame}{\end{frame}}}%
%BeginExpansion
\end{frame}%
%EndExpansion

%TCIMACRO{\TeXButton{BeginFrame}{\begin{frame}}}%
%BeginExpansion
\begin{frame}%
%EndExpansion

\QTR{frametitle}{Central Limit Theorem}

\begin{theorem}[The Central Limit Theorem]
Let $X_{1},...,X_{n}$ be i.i.d. with $\mathbb{E}\left( X_{1}\right) =\mu $
and $V\left( X_{1}\right) =\sigma ^{2}<\infty .$ Then, 
\begin{eqnarray*}
Z_{n} &=&\frac{\overline{X}_{n}-\mu }{\sqrt{V(\overline{X}_{n})}}=\frac{%
\sqrt{n}}{\sigma }(\overline{X}_{n}-\mu )=\frac{1}{\sigma \sqrt{n}}%
\sum_{i=1}^{n}(X_{i}-\mu )\rightarrow _{d}Z\text{ } \\
\text{\ where }Z &\sim &\mathcal{N}(0,1)
\end{eqnarray*}

In other words, 
\begin{equation*}
\lim_{n\rightarrow \infty }F_{Z_{n}}=\Phi (z)=\int_{-\infty }^{z}(2\pi
)^{-1/2}\exp (-x^{2}/2)dx.
\end{equation*}
\end{theorem}

%TCIMACRO{\TeXButton{EndFrame}{\end{frame}}}%
%BeginExpansion
\end{frame}%
%EndExpansion

%TCIMACRO{\TeXButton{BeginFrame}{\begin{frame}}}%
%BeginExpansion
\begin{frame}%
%EndExpansion

\QTR{frametitle}{Reflection }

\begin{itemize}
\item The LLN tells us that the distribution of $\overline{X}_{n}$ piles up
around $\mu .$

\item The CLT\ says that probability of $\frac{\sqrt{n}\left( \bar{X}%
_{n}-\mu \right) }{\sigma }$ converges to a $N\left( 0,1\right) $\ r.v.

\item This is a remarkable result, since nothing was assumed about the
distribution of the i.i.d. sequence $\left\{ X_{1},..,X_{n}\right\} $, apart
from the existence of their mean and variance; so the $X^{\prime }s$ could
be generated by any distribution.
\end{itemize}

%TCIMACRO{\TeXButton{EndFrame}{\end{frame}}}%
%BeginExpansion
\end{frame}%
%EndExpansion

%TCIMACRO{\TeXButton{BeginFrame}{\begin{frame}}}%
%BeginExpansion
\begin{frame}%
%EndExpansion

\QTR{frametitle}{Multivariate Extensions}

\begin{itemize}
\item Let $\mathbf{y}_{1},...,\mathbf{y}_{n}$ be an i.i.d. sequence of $%
k\times 1$ vectors, let $\mu =\mathbb{E}[\mathbf{y}_{1}]$ be the $k\times 1$
vector of expectations.
\end{itemize}

\begin{theorem}[The Weak Law of Large Numbers]
If $\mathbf{y}_{1},...,\mathbf{y}_{n}$ is an i.i.d. sequence of $k\times 1$
vectors, and $\mathbb{E}\left\Vert \mathbf{y}_{1}\right\Vert <\infty ,$ then 
$\frac{1}{n}\dsum\nolimits_{i=1}^{n}\mathbf{y}_{i}\rightarrow _{p}\mu .$
Equivalently, $\left\Vert \frac{1}{n}\dsum\nolimits_{i=1}^{n}\mathbf{y}_{i}-%
\mathbb{E}\mathbf{y}_{1}\right\Vert \rightarrow _{p}0.$
\end{theorem}

\begin{itemize}
\item This extension implies that each element in the vector $\overline{%
\mathbf{y}}$ converges, this is referred to as pointwise convergence.
\end{itemize}

%TCIMACRO{\TeXButton{EndFrame}{\end{frame}}}%
%BeginExpansion
\end{frame}%
%EndExpansion

%TCIMACRO{\TeXButton{BeginFrame}{\begin{frame}}}%
%BeginExpansion
\begin{frame}%
%EndExpansion

\QTR{frametitle}{Multivariate Extensions}

\begin{theorem}[The Central Limit Theorem]
If $\mathbf{y}_{1},...,\mathbf{y}_{n}$ is an i.i.d. sequence of $k\times 1$
vectors, let $\mu =\mathbb{E}[\mathbf{y}_{1}]$ be the $k\times 1$ vector of
expectations and let $V$ be a $k\times k$ p.d. matrix $V\left( \mathbf{y}%
_{1}\right) =\mathbb{E}\left( \mathbf{y}_{1}-\mathbb{E}[\mathbf{y}%
_{1}]\right) \left( \mathbf{y}_{1}-\mathbb{E}[\mathbf{y}_{1}]\right)
^{\prime }$ with $\left\Vert V\right\Vert <\infty ,$ then%
\begin{equation*}
Z_{n}=\sqrt{n}\left( \overline{\mathbf{y}}_{n}-\mu \right) =\frac{1}{\sqrt{n}%
}\sum_{i=1}^{n}(\mathbf{y}_{i}-\mu )\rightarrow _{d}\mathcal{N}(0,V)
\end{equation*}
\end{theorem}

\begin{itemize}
\item Here the limit is a multivariate normal random vector
\end{itemize}

%TCIMACRO{\TeXButton{EndFrame}{\end{frame}}}%
%BeginExpansion
\end{frame}%
%EndExpansion

%TCIMACRO{\TeXButton{BeginFrame}{\begin{frame}}}%
%BeginExpansion
\begin{frame}%
%EndExpansion

\QTR{frametitle}{Multivariate Extensions}

\begin{definition}
A $k\times 1$ vector $\mathbf{y}$ is multivariate normal with mean $\mu \in 
\mathbb{R}^{k}$ and covariance matrix $V\in \mathbb{R}^{k\times k}$ with $V$
p.d definite, denoted $\mathbf{y\sim }\mathcal{N}\left( \mu ,V\right) \ $if
it has a joint density%
\begin{equation*}
f_{\mathbf{y}}=\left( 2\pi \right) ^{-k/2}\det \left( V\right)
^{-1/2}e^{-0.5\left( \mathbf{y-}\mu \right) ^{\prime }V^{-1}\left( \mathbf{y-%
}\mu \right) }.
\end{equation*}
\end{definition}

\begin{definition}
If $\mathbf{y\sim }\mathcal{N}\left( \mu ,V\right) $ and $\mathbf{x=}c+B%
\mathbf{y,}$ then $\mathbf{x\sim }\mathcal{N}\left( c+B\mu ,BVB^{\prime
}\right) .$
\end{definition}

%TCIMACRO{\TeXButton{EndFrame}{\end{frame}}}%
%BeginExpansion
\end{frame}%
%EndExpansion

%TCIMACRO{\TeXButton{BeginFrame}{\begin{frame}}}%
%BeginExpansion
\begin{frame}%
%EndExpansion

\QTR{frametitle}{The Delta Method}

\begin{itemize}
\item The delta method is just an application of Taylor expansion

\item It is often used in econometrics to compute the limit distribution of
standardised and centred quantities.
\end{itemize}

\begin{theorem}[the Delta method]
If $\sqrt{n}\left( \hat{\mu}-\mu \right) \rightarrow _{d}z$ and $g\left(
.\right) $ is a continuous and differentiable function in a neighbourhood of 
$\mu ,$ then 
\begin{equation*}
\sqrt{n}\left( g\left( \hat{\mu}\right) -g\left( \mu \right) \right)
\rightarrow _{d}Gz
\end{equation*}%
where $G=\frac{\partial \left( g\left( \mu \right) \right) }{\partial \mu }.$
\end{theorem}

%TCIMACRO{\TeXButton{EndFrame}{\end{frame}}}%
%BeginExpansion
\end{frame}%
%EndExpansion

\section{Asymptotic Properties of OLS}

%TCIMACRO{\TeXButton{BeginFrame}{\begin{frame}}}%
%BeginExpansion
\begin{frame}%
%EndExpansion

\QTR{frametitle}{Consistency}

\begin{definition}
An estimator $\hat{\theta}_{n}$ of $\theta $ is called \textbf{consistent}
if $\hat{\theta}_{n}\rightarrow _{p}\theta $ as the sample size $n$ tends to 
$\infty $.
\end{definition}

%TCIMACRO{\TeXButton{EndFrame}{\end{frame}}}%
%BeginExpansion
\end{frame}%
%EndExpansion

%TCIMACRO{\TeXButton{BeginFrame}{\begin{frame}}}%
%BeginExpansion
\begin{frame}%
%EndExpansion

\QTR{frametitle}{Consistency of OLS}

\begin{theorem}[Consistency]
Under Assumptions 1-4, $p\lim_{n\rightarrow \infty }\widehat{\beta }%
_{n}^{OLS}=\beta .$
\end{theorem}

%TCIMACRO{\TeXButton{EndFrame}{\end{frame}}}%
%BeginExpansion
\end{frame}%
%EndExpansion

%TCIMACRO{\TeXButton{BeginFrame}{\begin{frame}}}%
%BeginExpansion
\begin{frame}%
%EndExpansion

\QTR{frametitle}{Consistency of OLS}

\begin{proof}
$p\lim_{n\rightarrow \infty }\widehat{\beta }_{n}=\beta +p\lim_{n\rightarrow
\infty }\left[ \left( \frac{1}{n}\tsum\nolimits_{i=1}^{n}\mathbf{x}_{i}%
\mathbf{x}_{i}^{\prime }\right) ^{-1}\frac{1}{n}\tsum\nolimits_{i=1}^{n}%
\mathbf{x}_{i}\varepsilon _{i}\right] $%
\begin{eqnarray*}
&=&\beta +p\lim_{n\rightarrow \infty }\left( \frac{1}{n}\tsum%
\nolimits_{i=1}^{n}\mathbf{x}_{i}\mathbf{x}_{i}^{\prime }\right)
^{-1}p\lim_{n\rightarrow \infty }\left( \frac{1}{n}\tsum\nolimits_{i=1}^{n}%
\mathbf{x}_{i}\varepsilon _{i}\right) \text{ {\tiny by CMT}} \\
&=&\beta +\left( p\lim_{n\rightarrow \infty }\frac{1}{n}\tsum%
\nolimits_{i=1}^{n}\mathbf{x}_{i}\mathbf{x}_{i}^{\prime }\right)
^{-1}p\lim_{n\rightarrow \infty }\left( \frac{1}{n}\tsum\nolimits_{i=1}^{n}%
\mathbf{x}_{i}\varepsilon _{i}\right) \text{ {\tiny by CMT and Ass. 4}} \\
&=&\beta +Q^{-1}\mathbb{E[}\mathbf{x}_{i}\varepsilon _{i}]\text{ {\tiny by
CMT, WLLN and Ass. 4}} \\
&=&\beta \text{ \ }
\end{eqnarray*}%
since $\mathbb{E[}\mathbf{x}_{i}\varepsilon _{i}]=\mathbb{E[}\mathbf{x}_{i}%
\mathbb{E}\left( \varepsilon _{i}|x_{i}\right) ]=0$ by Assumption 3, note
that because of the i.i.d. assumptions, $\mathbb{E[}\mathbf{x}%
_{i}\varepsilon _{i}]=\mathbb{E[}\mathbf{x}_{1}\varepsilon _{1}].$
\end{proof}

\begin{itemize}
\item $\widehat{\beta }_{n}$ is a consistent estimator of $\beta .$
\end{itemize}

%TCIMACRO{\TeXButton{EndFrame}{\end{frame}}}%
%BeginExpansion
\end{frame}%
%EndExpansion

%TCIMACRO{\TeXButton{BeginFrame}{\begin{frame}}}%
%BeginExpansion
\begin{frame}%
%EndExpansion

\QTR{frametitle}{Asymptotic Normality}

\begin{theorem}[Asymptotic Normality]
Under Assumptions 1-4, $\sqrt{n}(\widehat{\beta }_{n}-\beta )\rightarrow _{d}%
\mathcal{N}(0,\sigma ^{2}Q^{-1}).$
\end{theorem}

%TCIMACRO{\TeXButton{EndFrame}{\end{frame}}}%
%BeginExpansion
\end{frame}%
%EndExpansion

%TCIMACRO{\TeXButton{BeginFrame}{\begin{frame}}}%
%BeginExpansion
\begin{frame}%
%EndExpansion

\begin{proof}
$\sqrt{n}(\widehat{\beta }_{n}-\beta )=\left( \frac{1}{n}\tsum%
\nolimits_{i=1}^{n}\mathbf{x}_{i}\mathbf{x}_{i}^{\prime }\right) ^{-1}\frac{1%
}{\sqrt{n}}\tsum\nolimits_{i=1}^{n}\mathbf{x}_{i}\varepsilon _{i}$%
\begin{gather}
\left( \frac{1}{n}\tsum\nolimits_{i=1}^{n}\mathbf{x}_{i}\mathbf{x}%
_{i}^{\prime }\right) ^{-1}\rightarrow _{p}Q^{-1}\text{ {\tiny by WLLN, CMT
and Ass. 4}}  \notag \\
\frac{1}{\sqrt{n}}(\tsum\nolimits_{i=1}^{n}\mathbf{x}_{i}\varepsilon _{i}-%
\mathbb{E[}\mathbf{x}_{1}\varepsilon _{1}])\rightarrow _{d}\mathcal{N}(0,var(%
\mathbf{x}_{1}\varepsilon _{1}))\text{ {\tiny by the CLT}}  \notag \\
\frac{1}{\sqrt{n}}(\tsum\nolimits_{i=1}^{n}\mathbf{x}_{i}\varepsilon
_{i})\rightarrow _{d}\mathcal{N}(0,\sigma ^{2}Q)  \notag \\
\Rightarrow \sqrt{n}(\widehat{\beta }_{n}-\beta )\rightarrow _{d}\mathcal{N}%
(0,\sigma ^{2}Q^{-1}QQ^{-1})\text{ {\tiny by Slutsky Theorem}}  \notag \\
\sqrt{n}(\widehat{\beta }_{n}-\beta )\rightarrow _{d}\mathcal{N}(0,\sigma
^{2}Q^{-1}).  \label{clt}
\end{gather}
\end{proof}

since $\mathbb{E[}\mathbf{x}_{i}\varepsilon _{i}]=\mathbb{E[}\mathbf{x}%
_{1}\varepsilon _{1}]=0$ and $var(\mathbf{x}_{i}\varepsilon _{i})=var(%
\mathbf{x}_{1}\varepsilon _{1})=$ $\mathbb{E[}\mathbf{x}_{1}\mathbb{E}\left[
\varepsilon _{1}^{2}\right] \mathbf{x}_{1}^{\prime }]=\sigma ^{2}Q.$

%TCIMACRO{\TeXButton{EndFrame}{\end{frame}}}%
%BeginExpansion
\end{frame}%
%EndExpansion

%TCIMACRO{\TeXButton{BeginFrame}{\begin{frame}}}%
%BeginExpansion
\begin{frame}%
%EndExpansion

\QTR{frametitle}{Remarks}

\begin{itemize}
\item The asymptotic normality result tells us that for large enough $n$, we
can use the normal distribution for inference

\item So we can construct confidence intervals for $\beta $ and conduct
hypothesis testing

\item We obtained this result without assuming anything on the distribution
of $\varepsilon $

\item In fact, all results so far (Gauss Markov, unbiasedness, etc) were
derived \textit{without} assuming a particular distribution for $\varepsilon
;$ in this sense, OLS is a semi-parametric (or partial information) procedure

\item Matching moment estimators are all partial information methods

\item Maximum likelihood and Bayesian methods are full information, in the
sense that they require distributional assumptions
\end{itemize}

%TCIMACRO{\TeXButton{EndFrame}{\end{frame}}}%
%BeginExpansion
\end{frame}%
%EndExpansion

%TCIMACRO{\TeXButton{BeginFrame}{\begin{frame}}}%
%BeginExpansion
\begin{frame}%
%EndExpansion

\QTR{frametitle}{Remarks}

\begin{itemize}
\item If we assumed in Assumption 3 in addition that $\varepsilon |X\sim 
\mathcal{N}(0,\sigma ^{2}I_{n}),$ then the result in equation (\ref{clt})
holds without evoking the CLT; and crucially $\widehat{\beta }_{n}\sim 
\mathcal{N}$ even for small $n.$

\item We will revisit the Gaussian linear regression model when we study
maximum likelihood estimation, and we will state some small sample testing
procedures such as the $t$ and $F$ test

\item For now, we will rely on the large sample result

\item The advantage: it applies to any error distribution, the disadvantage:
it is an approximation that only works when $n$ is large
\end{itemize}

%TCIMACRO{\TeXButton{EndFrame}{\end{frame}}}%
%BeginExpansion
\end{frame}%
%EndExpansion

\section{Reading}

%TCIMACRO{\TeXButton{BeginFrame}{\begin{frame}}}%
%BeginExpansion
\begin{frame}%
%EndExpansion

\QTR{frametitle}{Reading}

\begin{itemize}
\item Relevant chapters from the textbook: 6 and 7.
\end{itemize}

%TCIMACRO{\TeXButton{EndFrame}{\end{frame}}}%
%BeginExpansion
\end{frame}%
%EndExpansion

\end{document}
